\chapter{المقاطعات وبرامج تشغيل الأجهزة}
\label{CH:INTERRUPT}

\textbf{برنامج تشغيل الجهاز} (Device Driver) هو قطعة البرمجية داخل النواة التي تدير وتدعم التفاعل مع جهاز عتادي معين (مثل القرص الصلب، بطاقة الشبكة، أو لوحة المفاتيح).
تتصل برامج تشغيل الأجهزة بالعتاد عبر السجلات العتادية المربوطة بالذاكرة (Memory-Mapped I/O) وتتلقى الإشعارات غير المتزامنة عبر \textbf{المقاطعات} (Interrupts).

يقدم هذا الفصل طريقة عمل المقاطعات العتادية، وإدارة الإدخال والإخراج عبر شاشة المنصة UART، ومقاطعات المؤقت الزمني في \texttt{xv6}.

\section{الكود البرمجي: مدخلات المنصة (UART Input)}
\label{sec:code_console_input}

يتواصل \texttt{xv6} مع المستخدم عبر جهاز منفذ تسلسلي يُدعى \textbf{UART 16550}.
عندما يكتب المستخدم حرفًا على لوحة المفاتيح:
\begin{enumerate}
  \item يولد عتاد UART إشارة مقاطعة متصلة بمجمع المقاطعات PLIC (Platform-Level Interrupt Controller).
  \item يرسل PLIC المقاطعة إلى إحدى أنوية المعالج.
  \item تنقل النواة السيطرة إلى دالة المصايد \texttt{devintr()} في \texttt{kernel/trap.c}.
  \item تتعرف \texttt{devintr()} على رقم المقاطعة القادمة من PLIC وتستدعي \texttt{uartintr()} المحددة في \texttt{kernel/uart.c}.
  \item تقرأ \texttt{uartintr()} الأحرف المتاحة من سجل UART الخارجي وتمررها إلى \texttt{consoleintr()} في \texttt{kernel/console.c}.
  \item تخزن \texttt{consoleintr()} الأحرف في مفرغ حلقي (Circular Buffer) وتستدعي \texttt{wakeup()} لإيقاظ العمليات المنتظرة لإدخال سطر جديد عبر \texttt{read()}.
\end{enumerate}

\section{الكود البرمجي: مخرجات المنصة (UART Output)}
\label{sec:code_console_output}

عندما تقوم عملية بنسخ نصوص للطباعة على الشاشة عبر الاستدعاء \texttt{write()}:
\begin{enumerate}
  \item تتلقى النواة البيانات وتدعو \texttt{uartputc()}.
  \item تُضاف البيانات إلى المخزن المؤقت للإخراج المكتبي في \texttt{kernel/uart.c}.
  \item إذا كان جهاز UART جاهزًا لإرسال البايتات، ينقل المحرك البيانات إلى سجل الإرسال \texttt{THR}.
  \item عند اكتمال إرسال كل بايت، يطلق UART مقاطعة إكتمال الإرسال، فتستجيب النواة بإرسال البايت التالي في المخزن المؤقت.
\end{enumerate}

\section{التزامن في برامج تشغيل الأجهزة}
\label{sec:concurrency_drivers}

تفرض برامج تشغيل الأجهزة التعامل مع مستويين من التزامن:
\begin{itemize}
  \item \textbf{التزامن بين النواة والجهاز}: يعالج العتاد الطلبات بطريقة غير متزامنة، مما يتطلب استخدام الأقفال (\texttt{spinlock}) لتعليق العمليات حتى يكتمل نقل البيانات عبر \texttt{sleep()} و \texttt{wakeup()}.
  \item \textbf{التزامن بين المقاطعات والأكواد النواتية}: لتجنب حدوث الجمود إذا وقُعت مقاطعة أثناء حيازة قفل معين، تقوم النواة بتعطيل المقاطعات على النواة الحالية فور حيازة قفل التناوب عبر \texttt{push\_off()}، وتفعلها بعد تحرير القفل عبر \texttt{pop\_off()}.
\end{itemize}

\section{مقاطعات المؤقت الزمني}
\label{sec:timer_interrupts}

يحتاج نظام التشغيل إلى مقاطعة دورية من المؤقت العتادي للتحكم في الجدولة الزمنية وتطبيق التعددية المشاركة (Preemptive Multitasking).

في RISC-V:
\begin{enumerate}
  \item تصدر مقاطعات المؤقت العتادية من مجمع المقاطعات المحلي \textbf{CLINT} وتُعالج أولاً في وضع الآلة (M-mode) داخل البرمجيات الثابتة.
  \item يقوم كود وضع الآلة (\texttt{kernel/start.c}) بإعداد المؤقت الدوري عبر سجل \texttt{mtimecmp}.
  \item تحول البرمجية المقاطعة إلى مقاطعة برمجية في وضع المشرف (Supervisor Software Interrupt).
  \item تتلقى النواة المقاطعة في \texttt{clockintr()} وتدعو \texttt{yield()} لإجبار العملية الحالية على التخلي عن وحدة المعالجة والسماح لعملية أخرى بالتنفيذ.
\end{enumerate}

\section{العالم الحقيقي}
\label{sec:real_world_interrupt}

في الأنظمة التجارية الحديثة، تتم أداء الإدخال والإخراج المتقدم للملفات والشبكات عبر تقنية \textbf{الوصول المباشر للذاكرة} (Direct Memory Access - DMA). تسمح تقنية DMA للأجهزة بنقل كتل كاملة من البيانات مباشرة من وإلى الذاكرة الفيزيائية دون إشغال وحدة المعالجة المركزية، حيث يقتصر دور المقاطعة على إعلام النواة بإنهاء الكتلة فقط.

\section{تمارين}
\label{sec:exercises_interrupt}

\begin{enumerate}
  \item ما السبب وراء تحذير استدعاء الأقفال داخل معالجات المقاطعات دون تعطيل المقاطعات أولاً؟
  \item أضف دعماً لجهاز إدخال لوحة مفاتيح حقيقي أو افتراضي إضافي في \texttt{xv6}.
\end{enumerate}
